\documentclass[12pt,a4paper]{article}
\usepackage[utf8]{inputenc}
\usepackage{amsmath}
\usepackage{amsfonts}
\usepackage{amssymb}
\usepackage{graphicx}
\usepackage{hyperref}
\usepackage{geometry}
\geometry{margin=1in}

\title{Connectiva: An Autonomous Telecommunications Optimization Platform}
\author{Team Connectiva \\ ITU UMC Hackathon 2026}
\date{\today}

\begin{document}

\maketitle

\section{Introduction}
Connectiva is a production-grade optimization engine designed to bridge the digital divide in Bangladesh and beyond. The platform leverages Machine Learning (Extra Trees Classifier) and Control Theory (PID) to simulate policy interventions and predict their multi-year socio-economic impact.

\section{Mathematical Model}

\subsection{Indicator Weighting and Feature Importance}
The platform utilizes a feature importance matrix derived from 25 years of ITU global data. The contribution of each indicator $x_i$ to the final connectivity score $S$ is defined by:
\begin{equation}
S = \sum_{i=1}^{n} w_i \cdot x_i \cdot \alpha_i
\end{equation}
where $w_i$ is the normalized MVT weight and $\alpha_i$ is the user-defined policy multiplier (ranging from 0.0 to 2.0).

\subsection{ConnectivaNet v4: Extra Trees Classifier}
For classification of district health (Green/Yellow/Red), we implement an Extra Trees Classifier. The split criterion is Gini Impurity:
\begin{equation}
G = 1 - \sum_{i=1}^{c} p_i^2
\end{equation}
This model provides higher variance reduction compared to standard Random Forests, making it ideal for the volatile geopolitical data seen in merging markets.

\subsection{Multi-Year Growth Simulation (PID)}
To simulate the roadmap for a district reaching its target $T$, we use a Proportional-Integral-Derivative (PID) controller. The annual growth adjustment $\Delta C_t$ is calculated as:
\begin{equation}
\Delta C_t = K_p e(t) + K_i \int e(t) dt + K_d \frac{de(t)}{dt}
\end{equation}
where $e(t) = T - C_{t-1}$. In our implementation, we use $K_p=0.6$, $K_i=0.1$, and $K_d=0.05$.

\section{Project Iterations}

\subsection{Iteration 1: Raw Data Ingestion}
The project began with the ingestion of raw BTRC (subscriber) and NTTN (fiber) data. Initial analysis showed a heavy bias toward Dhaka and Chittagong.

\subsection{Iteration 2: 25-Year Global Benchmarking}
We integrated ITU data spanning 1998-2023 to create global benchmarks. This allowed us to compare Bangladesh's "Gender Gap in Internet Usage" and "Fixed vs Mobile Broadband" growth against global leaders (South Korea, Singapore).

\subsection{Iteration 3: PID Optimization Engine}
The implementation of the PID controller allowed for realistic multi-year roadmaps, moving away from linear projections to non-linear, saturation-aware growth models.

\subsection{Iteration 4: Autonomous 2-Input Goal Logic}
The final iteration introduced a bidirectional dependency system where setting any two of \{Target, Timeframe, Budget\} automatically derives the most efficient value for the third, ensuring economic feasibility.

\section{Feature Analysis}
Topographically, Connectiva identifies "Priority Districts" using a 70:30 heuristic (70\% lowest Yellow, 30\% lowest Red), ensuring that both near-target and far-target regions receive balanced attention.

\section{Conclusion}
Connectiva provides a scalable, mathematics-first approach to national infrastructure planning. By combining ML-driven insights with real-time policy simulations, it empowers regulators to make data-backed decisions that maximize social ROI.

\end{document}
